% TOP_PROBLEM: {"id": 3401, "title": "Discrete Logarithm Problem", "category_id": 15, "category": {"id": 15, "name": "computer_science", "display_name": "Computer Science", "description": "Computational complexity, algorithms, and theoretical CS.", "slug": "computer-science", "order_index": 15, "created_at": "2026-07-31T15:26:25.671Z"}, "status": "open", "problem_number": "OPG-2150", "set_id": 12}
% FIRST_POSTED: 2026-10-02
% ATTEMPT_STATUS: unresolved
% MODEL: GPT 6 Astra Ultra
% UnsolvedMath ID 3401; source catalogue code OPG-2150
% !TeX program = lualatex
% Research attempt dated 2026-10-02; canonical statement preserved.
\documentclass[11pt,a4paper]{article}
\usepackage[margin=25mm]{geometry}
\usepackage{fontspec}
\setmainfont{lmroman10-regular.otf}[BoldFont=lmroman10-bold.otf,
 ItalicFont=lmroman10-italic.otf,BoldItalicFont=lmroman10-bolditalic.otf]
\usepackage{amsmath,amssymb,mathtools}
\usepackage{unicode-math}
\setmathfont{latinmodern-math.otf}
\AtBeginDocument{\let\setminus\smallsetminus}
\usepackage{xurl,hyperref}
\hypersetup{colorlinks=true,urlcolor=blue}
\setcounter{secnumdepth}{0}
\setlength{\parindent}{0pt}
\setlength{\parskip}{5pt}
\setlength{\emergencystretch}{3em}
\begin{document}
\section*{Discrete Logarithm Problem}\label{3401}
{\small UnsolvedMath ID 3401 · OPG-2150 · Computer Science · OpenGarden}
\subsection{Definitions and mathematical statement}
Let $p$ be prime and let $g,h$ be nonzero residues modulo $p$, given
by their least positive integer representatives. Assume $g\ne1$
and promise $h\in\langle g\rangle$, the cyclic subgroup generated
by $g$. A discrete logarithm of $h$ to base $g$ is an integer $a$
such that $g^a\equiv h\pmod p$. If $r=\operatorname{ord}_p(g)$,
all solutions are one residue class modulo $r$; an output
$0\le a<p-1$ suffices.

The conjecture, in the classical deterministic interpretation, is
that no uniform algorithm solves every promised input in time
polynomial in $\log p$. Input and output use binary encoding, and
time is measured in bit operations. The subgroup promise fixes a
missing condition in the short source: for an arbitrary pair $g,h$
a logarithm need not exist. The question includes the case that $g$
generates the whole multiplicative group.

The target is worst-case complexity over unbounded primes, not
hardness for a fixed modulus. It does not assert that an individual
logarithm cannot be found by exhaustive search. Quantum algorithms,
randomized algorithms, and distributional security claims use
different computational requirements and must be distinguished from
this precise classical deterministic assertion.
\subsection{Short English statement}
Is there no classical deterministic algorithm that finds every prime-field discrete logarithm in time polynomial in the bit length of the modulus?
\subsection{Research attempt}
\paragraph{Scope and method.}
This research attempt by GPT-6 Astra Ultra is dated 2 October 2026.
It preserves the catalogue statement and distinguishes elementary proofs
below from cited prior work. No novelty claim or formal proof-assistant
verification is made. The final paragraph states the precise remaining scope.
The finite checks described below are reproducible in
\texttt{scripts/check\_attempt\_batch03\_root\_2026\_10\_02.py}.
\paragraph{Bit complexity and the direction of the question.}
The input encodes $p$ in binary, so polynomial in $p$ is
exponential in its bit length. The catalogue asks for a
classical worst-case complexity conclusion over all valid
inputs. We give an unconditional general algorithm with
exponential bit complexity, and a polynomial algorithm for
a restricted family. Neither proves a polynomial algorithm
for every prime nor rules one out. The prime-power subgroup
method is a special case of Pohlig--Hellman [R1], whose
provenance matters: it is not a new discrete-logarithm
algorithm claimed by this attempt.

\paragraph{A general algorithm without knowing the order.}
Put $N=p-1$ and $m=\lceil\sqrt N\rceil$. Store the residues
$g^j$ for $0\le j<m$, together with an exponent for each
distinct residue. Compute $b=g^{-m}$ modulo $p$ using
the extended Euclidean algorithm for inversion. For
$i=0,\ldots,m-1$, test whether $hb^i$ occurs in the table.
A match with $g^j$ gives $g^{im+j}=h$. Reduce $im+j$
modulo $N$ to get an allowed output and verify it by
modular exponentiation.
Fermat's theorem gives $g^N=1$, and the promise on $h$
provides some exponent $0\le a<N$. Write $a=im+j$
with $0\le j<m$. Because $N\le m^2$, its index satisfies
$i<m$, and that iteration finds a matching residue.
If the table has retained a different exponent for that
same residue, it still yields a correct solution.
Thus repeated powers when $g$ is not a generator do
not invalidate the method. A sorted table gives a
deterministic implementation with $O(m\log m)$ residue
comparisons and $O(m)$ modular multiplications, apart
from polynomial-size preliminary arithmetic. This is
$p^{1/2}$ times polynomial factors in $\log p$, which
is still exponential in the input length.

\paragraph{An exact polynomial subfamily.}
Suppose $p-1=2^s$. Repeatedly square $g$ until obtaining
one; the first exponent is $r=2^t$, the exact order of
$g$, with $1\le t\le s$. The promise excludes $g=1$.
The element $g^{2^{t-1}}$ is the unique element of order
two in $\mathbb F_p^*$, hence equals $-1$.
Recover a logarithm modulo $2^t$ one bit at a time.
If $a_j$ contains the already recovered bottom $j$
bits, form
\[
 z_j=(h g^{-a_j})^{2^{t-j-1}}.
\]
Writing $a-a_j=2^j b$, correctness of the previous
bits gives $z_j=(-1)^b$. Thus $z_j=1$ means that
the next bit is zero, and $z_j=-1$ means it is one.
Set $a_{j+1}=a_j$ or $a_j+2^j$ accordingly, starting
with $a_0=0$. After $t$ steps the resulting $a_t$
lies in $[0,r)$ and satisfies $g^{a_t}=h$.
There are at most $s$ modular exponentiations of
polynomial bit complexity, so the whole restricted
algorithm is polynomial in $\log p$. No assumption
that infinitely many such primes exist is needed.

\paragraph{Where the general barrier remains.}
The bit extraction works because successive quotients
of the cyclic subgroup have order two. Replacing a
large prime-order quotient by a two-way sign test is
unjustified. More generally, a fast algorithm when
the group order has small prime factors does not
cover an arbitrary prime $p$. Conversely, a lower
bound for algorithms that access only a generic
group oracle would not automatically cover algorithms
that exploit the binary representation and field
arithmetic of these inputs. Verifying a proposed
logarithm efficiently also provides neither a fast
search algorithm nor a hardness proof.

\paragraph{Verification and final status.}
The root script checks the general algorithm on
all nonidentity bases and generated targets for
small primes, including nongenerators and $h=1$.
It checks every base and subgroup target for
$p=3,5,17,257$ in the bit-extraction algorithm.
These are exact modular computations. The proofs
establish their specified domains, but the unrestricted
classical polynomial-time question remains unresolved
by this attempt.

\subsection{Sources}
\begin{itemize}
\item[S1] ULAM.AI and the UnsolvedMath Contributors, supplied problems.json, record 3401 (OPG-2150), OpenGarden collection. Catalogue identity and source transcription; CC BY 4.0.
\url{https://huggingface.co/datasets/ulamai/UnsolvedMath}.
\item[S2] Open Problem Garden, Discrete Logarithm Problem. Original problem and discussion; the mathematical formulation below is expanded from the supplied source record.
\url{http://www.openproblemgarden.org/op/discrete_logarithm_problem}.

\item[R1] S. C. Pohlig and M. E. Hellman, An improved algorithm for
computing logarithms over GF(p) and its cryptographic significance,
IEEE Trans. Information Theory 24 (1978), 106--110.
\url{https://doi.org/10.1109/TIT.1978.1055817}.
Author-hosted paper: \url{https://ee.stanford.edu/~hellman/publications/28.pdf}.

\end{itemize}
\end{document}
